We perform learning rate hyperparameter tuning for our evaluation models and consider using either random weights or pretrained weights. Selecting the best possible evaluation model is crucial as it can greatly affect the evaluation of our augmentation techniques. An underperforming model may produce incorrect predictions, leading to inaccurate results and unreliable conclusions.

\subsubsection{YOLOv5s}
Two YOLOv5 models were trained without augmentation for 100 epochs: one with pretrained weights from the COCO dataset, and one without pretrained weights. Results in Figure \ref{fig:baseline_yolo}a) demonstrate that the pretrained model outperforms the non-pretrained model. This suggests that the general features learned during pretraining, such as edges, corners, and textures, are still useful for detection in our dataset even though the COCO is not specific to aerial imagery \cite{cocodataset}.

To further optimize performance, a hyperparameter tuning experiment was conducted to determine the optimal learning rate. The results, shown in Figure \ref{fig:baseline_yolo}b), indicate that the best learning rate was 0.01, which achieved an F1-Score top-3 of 0.1773. To summarize, we have determined that YOLOv5s with pretrained weights and a learning rate of 0.01 is the baseline for the deteciton task.

\begin{figure}[!h]
\centerline{\subfigure[With and without pretrained weights]{\includegraphics[width=12cm]{src/Results/baseline/base_detec_pre.png}}}
\centerline{\subfigure[Learning rate tuning]{\includegraphics[width=12cm]{src/Results/baseline/detec_lr.png}}}
\caption{Illustration of (a) YOLOv5s with and without pretrained weights and (b) on different learning rate}
\label{fig:baseline_yolo}
\end{figure}

% \begin{figure}[!h]
% 	 \centerline{\includegraphics[width=0.7\textwidth]{src/Results/baseline/base_detec_pre.png}}
% 	 {\caption{mean average precision of Yolov5 Lite with and without pretrained weights}\label{fig:base_detec_pre}}
% \end{figure}

% \begin{figure}[!h]
% 	 \centerline{\includegraphics[width=0.7\textwidth]{src/Results/baseline/detec_lr.png}}
% 	 {\caption{Learning Rate Hyperparatmeter tuning for Yolov5}\label{fig:detec_lr}}
% \end{figure}
\subsubsection{EfficientNet}

The EfficientNet was trained for 200 epochs using both random-initialized weights and pretrained weights from the Imagenet dataset, without any augmentation. The results, as illustrated in figure \ref{fig:baseline_effi}a), demonstrate that the pretrained model outperformed the model without pretrained weights, suggesting that learning transferability exists even when the datasets lack significant common features.

After conducting hyperparameter tuning within a logarithmic search space between $10^{-4}$ and $10^{-8}$, we determined that the optimal learning rate for EfficientNet version B0 with pretrained weights was $10^{-5}$, as demonstrated in figure \ref{fig:baseline_effi}b). This yielded an F1-Score top-3 of 0.3434. Thus, we established the baseline for the classification task as using EfficientNet version B0 with pretrained weights and a learning rate of $10^{-5}$.

% \begin{figure}[!h]
% 	 \centerline{\includegraphics[width=1\textwidth]{src/Results/baseline/base-class-pre.png}}
% 	 {\caption{F1-Score top-3 of EfficienNet B0 with and without pretrained weights}\label{fig:base-class-pre}}
% \end{figure}

% \begin{figure}[!h]
% 	 \centerline{\includegraphics[width=1\textwidth]{src/Results/baseline/class_lr.png}}
% 	 {\caption{Learning Rate Hyperparatmeter tuning for EfficientNet}\label{fig:class_lr}}
% \end{figure}

\begin{figure}[!ht]
\centerline{\subfigure[With and without pretrained weights]{\includegraphics[width=12cm]{src/Results/baseline/base-class-pre.png}}}
\centerline{\subfigure[Learning rate tuning]{\includegraphics[width=12cm]{src/Results/baseline/class_lr.png}}}
\caption{Illustration of EfficientNet version B0 (a) with and without pretrained weights and (b) on different learning rates}
\label{fig:baseline_effi}
\end{figure}
